% GENERATED FILE. Edit canonical lessons, then run npm run generate:books.
% canonical-source-hash: fnv1a64:38c53ae785321f35
% canonical-lessons: GU-C468-juthu, GU-C468-gappu, GU-C468-dambh, GU-C468-sekhi, GU-C468-bholavvu

\chapter{Lies and Boasts}
\label{ch:gu-lies-and-boasts}
\hlchaptermodality{468}

\begin{chapteropening}
\textbf{By the end of this chapter:} \emph{I can say when someone is not telling the truth.}

Complete the last lesson of chapter 468: I can say when someone is not telling the truth.
\end{chapteropening}

\section[\texorpdfstring{\gu{જૂઠું} (jūṭhũ)}{જૂઠું (jūṭhũ)}]{\gu{જૂઠું} (jūṭhũ) — untruthful; lying}

\label{lesson:GU-C468-juthu}



\begin{quote}
Before the new one: say the Gujarati for moral character, then the Gujarati for to be ashamed.
\end{quote}

\subsection*{You'll want to know: \gu{જૂઠું}}
\textbf{\gu{જૂઠું}} — \emph{jūṭhũ} — \textquotedblleft{}untruthful; lying\textquotedblright{}.

\begin{quote}
\gu{જૂઠું} \gu{બોલવું} \gu{સારું} \gu{નથી}.
\end{quote}

\begin{quote}
\emph{jūṭhũ bolvũ sārũ nathī.}
\end{quote}

\begin{quote}
Telling lies is not good.
\end{quote}

\begin{cousinweb}
\textbf{\gu{જૂઠું}} (\emph{jūṭhũ}) goes back to a reconstructed \emph{*jhūṭṭha-}, \textquotedblleft{}false\textquotedblright{}; CDIAL lists the Gujarati word, \textquotedblleft{}false\textquotedblright{} (CDIAL 5407).
\end{cousinweb}

\begin{grammarlens}[title={What you've built}]
One more word for describing what you lived and felt.
\end{grammarlens}

\subsection*{Your turn}
\begin{itemize}
\raggedright
  \item \emph{Say it:} \emph{jūṭhũ}
  \item \emph{Say it:} \emph{jūṭhũ}, once more
  \item \emph{Say it:} \emph{jūṭhũ bolvũ sārũ nathī.} — Telling lies is not good.
  \item \emph{Recall:} say the Gujarati for moral character, then the Gujarati for to be ashamed, then say \emph{jūṭhũ} again
\end{itemize}

\begin{tcolorbox}[breakable,colback=teal!4,colframe=teal!35!black,title={Before you move on}]
What does \gu{જૂઠું} mean? (\textquotedblleft{}Untruthful; lying\textquotedblright{}.) Say it once more.
\end{tcolorbox}
\section[\texorpdfstring{\gu{ગપ્પું} (gappũ)}{ગપ્પું (gappũ)}]{\gu{ગપ્પું} (gappũ) — a tall tale}

\label{lesson:GU-C468-gappu}



\begin{quote}
Before the new one: say the Gujarati for to be ashamed, then the Gujarati for untruthful.
\end{quote}

\subsection*{You'll want to know: \gu{ગપ્પું}}
\textbf{\gu{ગપ્પું}} — \emph{gappũ} — \textquotedblleft{}a tall tale\textquotedblright{}.

\begin{quote}
\gu{દાદાનાં} \gu{ગપ્પાં} \gu{સાંભળીને} \gu{બધા} \gu{હસ્યા}.
\end{quote}

\begin{quote}
\emph{dādānā̃ gappā̃ sā̃bhḷīne badhā hasyā.}
\end{quote}

\begin{quote}
Everyone laughed at my grandfather's tall tales.
\end{quote}

A neuter noun.

\begin{cousinweb}
\textbf{\gu{ગપ્પું}} (\emph{gappũ}) goes back to a reconstructed \emph{*gappa-}, \textquotedblleft{}talk\textquotedblright{}, probably an imitative word; CDIAL lists the Gujarati \emph{gap}, \textquotedblleft{}tittle-tattle\textquotedblright{} (CDIAL 4022).
\end{cousinweb}

\begin{grammarlens}[title={What you've built}]
One more word for describing what you lived and felt.
\end{grammarlens}

\subsection*{Your turn}
\begin{itemize}
\raggedright
  \item \emph{Say it:} \emph{gappũ}
  \item \emph{Say it:} \emph{gappũ}, once more
  \item \emph{Say it:} \emph{dādānā̃ gappā̃ sā̃bhḷīne badhā hasyā.} — Everyone laughed at my grandfather's tall tales.
  \item \emph{Recall:} say the Gujarati for to be ashamed, then the Gujarati for untruthful, then say \emph{gappũ} again
\end{itemize}

\begin{tcolorbox}[breakable,colback=teal!4,colframe=teal!35!black,title={Before you move on}]
What does \gu{ગપ્પું} mean? (\textquotedblleft{}A tall tale\textquotedblright{}.) Say it once more.
\end{tcolorbox}
\section[\texorpdfstring{\gu{દંભ} (dambh)}{દંભ (dambh)}]{\gu{દંભ} (dambh) — hypocrisy; pretence}

\label{lesson:GU-C468-dambh}



\begin{quote}
Before the new one: say the Gujarati for untruthful, then the Gujarati for a tall tale.
\end{quote}

\subsection*{You'll want to know: \gu{દંભ}}
\textbf{\gu{દંભ}} — \emph{dambh} — \textquotedblleft{}hypocrisy; pretence\textquotedblright{}.

\begin{quote}
\gu{તેનું} \gu{સ્મિત} \gu{ફક્ત} \gu{દંભ} \gu{હતું}.
\end{quote}

\begin{quote}
\emph{tenũ smit phakt dambh hatũ.}
\end{quote}

\begin{quote}
His smile was nothing but pretence.
\end{quote}

A masculine noun.

\begin{cousinweb}
\textbf{\gu{દંભ}} (\emph{dambh}) is a tatsama, Sanskrit \emph{dambha}, \textquotedblleft{}deceit, feigning, hypocrisy\textquotedblright{}, as Monier-Williams records it.
\end{cousinweb}

\begin{grammarlens}[title={What you've built}]
One more word for describing what you lived and felt.
\end{grammarlens}

\subsection*{Your turn}
\begin{itemize}
\raggedright
  \item \emph{Say it:} \emph{dambh}
  \item \emph{Say it:} \emph{dambh}, once more
  \item \emph{Say it:} \emph{tenũ smit phakt dambh hatũ.} — His smile was nothing but pretence.
  \item \emph{Recall:} say the Gujarati for untruthful, then the Gujarati for a tall tale, then say \emph{dambh} again
\end{itemize}

\begin{tcolorbox}[breakable,colback=teal!4,colframe=teal!35!black,title={Before you move on}]
What does \gu{દંભ} mean? (\textquotedblleft{}Hypocrisy; pretence\textquotedblright{}.) Say it once more.
\end{tcolorbox}
\section[\texorpdfstring{\gu{શેખી} (śekhī)}{શેખી (śekhī)}]{\gu{શેખી} (śekhī) — boasting}

\label{lesson:GU-C468-sekhi}



\begin{quote}
Before the new one: say the Gujarati for a tall tale, then the Gujarati for hypocrisy.
\end{quote}

\subsection*{You'll want to know: \gu{શેખી}}
\textbf{\gu{શેખી}} — \emph{śekhī} — \textquotedblleft{}boasting\textquotedblright{}.

\begin{quote}
\gu{અરુણ} \gu{પોતાની} \gu{નવી} \gu{ગાડીની} \gu{શેખી} \gu{મારે} \gu{છે}.
\end{quote}

\begin{quote}
\emph{Aruṇ potānī navī gāḍīnī śekhī māre che.}
\end{quote}

\begin{quote}
Arun boasts about his new car.
\end{quote}

A feminine noun.

\begin{cousinweb}
\textbf{\gu{શેખી}} (\emph{śekhī}) is usually traced to Persian \emph{šēxī}, \textquotedblleft{}pride, boasting\textquotedblright{}, built on \emph{šēx}, \textquotedblleft{}an elder, a chief\textquotedblright{}.
\end{cousinweb}

\begin{grammarlens}[title={What you've built}]
One more word for describing what you lived and felt.
\end{grammarlens}

\subsection*{Your turn}
\begin{itemize}
\raggedright
  \item \emph{Say it:} \emph{śekhī}
  \item \emph{Say it:} \emph{śekhī}, once more
  \item \emph{Say it:} \emph{Aruṇ potānī navī gāḍīnī śekhī māre che.} — Arun boasts about his new car.
  \item \emph{Recall:} say the Gujarati for a tall tale, then the Gujarati for hypocrisy, then say \emph{śekhī} again
\end{itemize}

\begin{tcolorbox}[breakable,colback=teal!4,colframe=teal!35!black,title={Before you move on}]
What does \gu{શેખી} mean? (\textquotedblleft{}Boasting\textquotedblright{}.) Say it once more.
\end{tcolorbox}
\section[\texorpdfstring{\gu{ભોળવવું} (bhoḷavvũ)}{ભોળવવું (bhoḷavvũ)}]{\gu{ભોળવવું} (bhoḷavvũ) — to beguile; to mislead}

\label{lesson:GU-C468-bholavvu}



\begin{quote}
Before the new one: say the Gujarati for hypocrisy, then the Gujarati for boasting.
\end{quote}

\subsection*{You'll want to know: \gu{ભોળવવું}}
\textbf{\gu{ભોળવવું}} — \emph{bhoḷavvũ} — \textquotedblleft{}to beguile; to mislead\textquotedblright{}.

\begin{quote}
\gu{લુચ્ચો} \gu{વેપારી} \gu{ગ્રાહકોને} \gu{ખોટા} \gu{વચનથી} \gu{ભોળવે} \gu{છે}.
\end{quote}

\begin{quote}
\emph{lucco vepārī grāhkone khoṭā vacanthī bhoḷve che.}
\end{quote}

\begin{quote}
The sly trader beguiles customers with false promises.
\end{quote}

Past, with \gu{તેણે} (he or she …): \gu{તેણે} \gu{બાળકને} \gu{ભોળવ્યું}, \emph{teṇe bāḷakne bhoḷavyũ}.

\begin{cousinweb}
\textbf{\gu{ભોળવવું}} (\emph{bhoḷavvũ}) is built on \textbf{\gu{ભોળું}} (\emph{bhoḷũ}), \textquotedblleft{}trusting\textquotedblright{}, which goes back to a reconstructed \emph{*bhulla-} (CDIAL 9539): to make someone trusting, and so to fool them.
\end{cousinweb}

\begin{grammarlens}[title={What you've built}]
One more word for describing what you lived and felt.
\end{grammarlens}

\subsection*{Your turn}
\begin{itemize}
\raggedright
  \item \emph{Say it:} \emph{bhoḷavvũ}
  \item \emph{Say it:} \emph{bhoḷavvũ}, once more
  \item \emph{Say it:} \emph{lucco vepārī grāhkone khoṭā vacanthī bhoḷve che.} — The sly trader beguiles customers with false promises.
  \item \emph{Recall:} say the Gujarati for hypocrisy, then the Gujarati for boasting, then say \emph{bhoḷavvũ} again
\end{itemize}

\begin{tcolorbox}[breakable,colback=teal!4,colframe=teal!35!black,title={Before you move on}]
What does \gu{ભોળવવું} mean? (\textquotedblleft{}To beguile; to mislead\textquotedblright{}.) Say it once more.
\end{tcolorbox}
